%% 第二章 无穷级数基础理论（约3000字）

\chapter{无穷级数基础理论}

本章系统回顾无穷级数的基本概念、收敛判别方法以及解析解与数值解的界定标准，为后续章节的深入研究奠定理论基础。

\section{无穷级数基本概念}

\subsection{常数项级数}

\begin{definition}
设 $\{a_n\}$ 是一个数列，称形式和
\begin{equation}
\sum_{n=1}^{\infty} a_n = a_1 + a_2 + a_3 + \cdots
\end{equation}
为\textbf{常数项无穷级数}（简称级数），其中 $a_n$ 称为级数的\textbf{通项}。
\end{definition}

级数的收敛性由其部分和序列决定，定义级数的第 $n$ 个部分和为
\begin{equation}
S_n = \sum_{k=1}^{n} a_k = a_1 + a_2 + \cdots + a_n
\end{equation}
若部分和序列 $\{S_n\}$ 收敛，即极限 $\displaystyle \lim_{n \to \infty} S_n = S$ 存在且有限，则称级数 $\displaystyle \sum_{n=1}^{\infty} a_n$ \textbf{收敛}，$S$ 称为级数的\textbf{和}；否则称级数\textbf{发散}。

级数收敛的一个必要条件是通项趋于零，即若 $\displaystyle \sum_{n=1}^{\infty} a_n$ 收敛，则 $\displaystyle \lim_{n \to \infty} a_n = 0$。但此条件并非充分条件，经典的反例是调和级数 $\displaystyle \sum_{n=1}^{\infty} \frac{1}{n}$，其通项趋于零但级数发散。

常见的常数项级数包括：
\begin{enumerate}[label=（\arabic*）]
  \item \textbf{等比级数}（几何级数）：$\displaystyle \sum_{n=0}^{\infty} ar^n$，当 $\abs{r} < 1$ 时收敛于 $\frac{a}{1-r}$；
  \item \textbf{$p$-级数}：$\displaystyle \sum_{n=1}^{\infty} \frac{1}{n^p}$，当 $p > 1$ 时收敛，$p \leq 1$ 时发散；
  \item \textbf{交错调和级数}：$\displaystyle \sum_{n=1}^{\infty} \frac{(-1)^{n+1}}{n}$，收敛于 $\ln 2$。
\end{enumerate}

\subsection{函数项级数}

\begin{definition}
设 $\{u_n(x)\}$ 是定义在区间 $I$ 上的函数列，称
\begin{equation}
\sum_{n=1}^{\infty} u_n(x) = u_1(x) + u_2(x) + u_3(x) + \cdots
\end{equation}
为定义在 $I$ 上的\textbf{函数项级数}。对于 $I$ 中的某个固定点 $x_0$，若常数项级数 $\displaystyle \sum_{n=1}^{\infty} u_n(x_0)$ 收敛，则称 $x_0$ 为该函数项级数的\textbf{收敛点}；所有收敛点的集合称为\textbf{收敛域}。
\end{definition}

函数项级数中最重要的两类是\textbf{幂级数}和\textbf{傅里叶级数}，分别介绍如下。

\textbf{幂级数}的一般形式为
\begin{equation}
\sum_{n=0}^{\infty} c_n (x - x_0)^n = c_0 + c_1(x - x_0) + c_2(x - x_0)^2 + \cdots
\end{equation}
其中 $c_n$ 为系数，$x_0$ 为展开中心。幂级数的收敛域具有对称性，存在\textbf{收敛半径} $R$（$0 \leq R \leq +\infty$），使得当 $\abs{x - x_0} < R$ 时级数绝对收敛，当 $\abs{x - x_0} > R$ 时级数发散。收敛半径可由以下公式计算（若极限存在）：
\begin{equation}
R = \lim_{n \to \infty} \left| \frac{c_n}{c_{n+1}} \right|
\end{equation}
或
\begin{equation}
R = \frac{1}{\limsup_{n \to \infty} \sqrt[n]{\abs{c_n}}}
\end{equation}

\textbf{傅里叶级数}将周期函数展开为三角函数的无穷和：
\begin{equation}
f(x) = \frac{a_0}{2} + \sum_{n=1}^{\infty} \left( a_n \cos nx + b_n \sin nx \right)
\end{equation}
其中傅里叶系数为
\begin{equation}
a_n = \frac{1}{\pi} \int_{-\pi}^{\pi} f(x) \cos nx \, \mathrm{d}x, \quad b_n = \frac{1}{\pi} \int_{-\pi}^{\pi} f(x) \sin nx \, \mathrm{d}x
\end{equation}

\section{收敛判别方法}

收敛判别法是无穷级数理论的核心工具，用于在无法直接计算部分和极限时判定级数的收敛性。

\subsection{正项级数的判别法}

正项级数是指每一项均非负的级数，即 $a_n \geq 0$ 对一切 $n$ 成立。正项级数的部分和序列单调递增，因此要么收敛要么趋向正无穷。

\textbf{比较判别法。}

\begin{theorem}[比较判别法]
设 $\displaystyle \sum a_n$ 和 $\displaystyle \sum b_n$ 均为正项级数，且存在正整数 $N$，使得当 $n \geq N$ 时 $a_n \leq b_n$，则：

（1）若 $\displaystyle \sum b_n$ 收敛，则 $\displaystyle \sum a_n$ 收敛；

（2）若 $\displaystyle \sum a_n$ 发散，则 $\displaystyle \sum b_n$ 发散。
\end{theorem}

比较判别法的核心思想是利用已知敛散性的级数作为"标尺"。在实际应用中，常以 $p$-级数和等比级数作为比较对象。例如，判别级数 $\displaystyle \sum_{n=1}^{\infty} \frac{1}{n^2 + 1}$ 的收敛性时，注意到 $\frac{1}{n^2 + 1} < \frac{1}{n^2}$，而 $\displaystyle \sum \frac{1}{n^2}$ 收敛（$p = 2 > 1$），故原级数收敛。

\textbf{比值判别法（达朗贝尔判别法）。}

\begin{theorem}[比值判别法]
设 $\displaystyle \sum a_n$ 为正项级数，若极限
\begin{equation}
\rho = \lim_{n \to \infty} \frac{a_{n+1}}{a_n}
\end{equation}
存在，则：当 $\rho < 1$ 时级数收敛；当 $\rho > 1$ 时级数发散；当 $\rho = 1$ 时判别法失效。
\end{theorem}

比值判别法特别适用于通项含有阶乘或指数因子的级数。例如，对级数 $\displaystyle \sum_{n=0}^{\infty} \frac{x^n}{n!}$，有 $\displaystyle \rho = \lim_{n \to \infty} \frac{\abs{x}}{n+1} = 0 < 1$，故该级数对一切实数 $x$ 绝对收敛，其和为 $e^x$。

\textbf{根值判别法（柯西判别法）。}

\begin{theorem}[根值判别法]
设 $\displaystyle \sum a_n$ 为正项级数，若极限
\begin{equation}
\rho = \lim_{n \to \infty} \sqrt[n]{a_n}
\end{equation}
存在，则：当 $\rho < 1$ 时级数收敛；当 $\rho > 1$ 时级数发散；当 $\rho = 1$ 时判别法失效。
\end{theorem}

根值判别法的适用范围比比值判别法更广，在通项含有 $n$ 次幂的结构时较为方便。例如，级数 $\displaystyle \sum_{n=1}^{\infty} \left(\frac{n}{2n+1}\right)^n$ 中，$\sqrt[n]{a_n} = \frac{n}{2n+1} \to \frac{1}{2} < 1$，故级数收敛。

\textbf{积分判别法（柯西积分判别法）。}

\begin{theorem}[积分判别法]
设 $f(x)$ 是定义在 $[1, +\infty)$ 上的非负单调递减函数，且 $a_n = f(n)$，则正项级数 $\displaystyle \sum_{n=1}^{\infty} a_n$ 与反常积分 $\displaystyle \int_1^{+\infty} f(x)\, \mathrm{d}x$ 同敛散：积分收敛则级数收敛，积分发散则级数发散。
\end{theorem}

积分判别法的一个重要推论是余项的积分估计：若正项级数 $\displaystyle \sum_{n=1}^{\infty} f(n)$ 收敛，则截断误差满足
\begin{equation} \label{eq:int_err}
R_N = \sum_{n=N+1}^{\infty} f(n) \leq \int_N^{+\infty} f(x)\, \mathrm{d}x.
\end{equation}
这一估计在后续 $p$-级数的数值求解与误差控制中将被频繁使用。例如，对 $p$-级数 $\displaystyle \sum_{n=1}^{\infty} 1/n^p$（$p > 1$），由公式\eqref{eq:int_err}可得截断误差上界 $R_N \leq N^{1-p}/(p-1)$。

\subsection{交错级数的判别法}

\textbf{莱布尼兹判别法。}

\begin{theorem}[莱布尼兹判别法]
设交错级数 $\displaystyle \sum_{n=1}^{\infty} (-1)^{n+1} b_n$（$b_n > 0$），若满足：

（1）$\{b_n\}$ 单调递减，即 $b_{n+1} \leq b_n$；

（2）$\displaystyle \lim_{n \to \infty} b_n = 0$，

则该交错级数收敛。
\end{theorem}

莱布尼兹判别法还给出了余项估计：截取前 $n$ 项后的误差不超过第 $n+1$ 项的绝对值，即 $\abs{S - S_n} \leq b_{n+1}$。这一性质在后续数值计算中将被使用，它为数值解的精度控制提供了直接依据。

交错调和级数 $\displaystyle \sum_{n=1}^{\infty} \frac{(-1)^{n+1}}{n}$ 是莱布尼兹判别法的经典应用实例，其收敛于 $\ln 2$。

\subsection{绝对收敛与条件收敛}

\begin{definition}
若级数 $\displaystyle \sum_{n=1}^{\infty} \abs{a_n}$ 收敛，则称 $\displaystyle \sum_{n=1}^{\infty} a_n$ \textbf{绝对收敛}；若 $\displaystyle \sum_{n=1}^{\infty} a_n$ 收敛但 $\displaystyle \sum_{n=1}^{\infty} \abs{a_n}$ 发散，则称 $\displaystyle \sum_{n=1}^{\infty} a_n$ \textbf{条件收敛}。
\end{definition}

绝对收敛是比条件收敛更强的性质：绝对收敛的级数必定收敛，反之不然。更重要的是，绝对收敛的级数在实数域中具有\textbf{无条件收敛}性——任意重排其各项后所得级数仍收敛且和不变。而条件收敛的级数则不然：根据黎曼重排定理，条件收敛级数的项可经重排使其收敛于任意预指定的值，甚至发散。

从本文研究角度而言，绝对收敛性对解析解的求取提供了更好的保障——绝对收敛级数可以进行逐项求导、逐项积分等运算而不改变收敛性，这为解析方法的实施提供了理论基础。

\section{解析解与数值解的界定}

在无穷级数的求解实践中，解的存在形式并非单一的。本节明确解析解与数值解的概念，为后续的谱系化研究提供分类依据。

\subsection{解析解的定义}

\begin{definition}
若级数 $\displaystyle \sum_{n=1}^{\infty} a_n$ 的和 $S$ 可以用\textbf{有限次}基本运算（加、减、乘、除）和初等函数（幂函数、指数函数、对数函数、三角函数及其复合）的\textbf{闭合表达式}精确表示，则称 $S$ 为该级数的\textbf{解析解}。
\end{definition}

典型的解析解实例包括：
\begin{enumerate}[label=（\arabic*）]
  \item 等比级数：$\displaystyle \sum_{n=0}^{\infty} r^n = \frac{1}{1-r}$（$\abs{r} < 1$）；
  \item 巴塞尔问题：$\displaystyle \sum_{n=1}^{\infty} \frac{1}{n^2} = \frac{\pi^2}{6}$；
  \item 交错调和级数：$\displaystyle \sum_{n=1}^{\infty} \frac{(-1)^{n+1}}{n} = \ln 2$。
\end{enumerate}

解析解的优势在于其精确性和普适性——它不受计算精度的限制，能够揭示级数和与参数之间的函数关系。

\subsection{数值解的定义}

\begin{definition}
若级数 $\displaystyle \sum_{n=1}^{\infty} a_n$ 的和 $S$ 无法用闭合表达式精确表示，通过截取有限项部分和 $S_n$ 或其他数值方法所获得的\textbf{近似值}称为该级数的\textbf{数值解}。即便某级数理论上存在解析解，在实际数值计算中若采用截断或迭代方法求值，所得结果同样属于数值近似的范畴。
\end{definition}

数值解通常表现为一个带有误差界的近似数值。例如，级数 $\displaystyle \sum_{n=1}^{\infty} \frac{1}{n^3}$（阿佩里常数相关级数）虽然已知收敛，但其和 $\zeta(3) \approx 1.2020569\ldots$ 至今没有已知的闭合表达式，只能通过数值方法逐步逼近。

\subsection{两者的比较}

解析解与数值解的核心区别在于表达方式和精确程度，两者各具优势与局限，具体比较见表~\ref{tab:compare}。

\begin{table}[htbp]
\centering
\caption{解析解与数值解的比较}
\label{tab:compare}
\begin{tabularx}{\textwidth}{lXX}
\toprule
\textbf{比较维度} & \textbf{解析解} & \textbf{数值解} \\
\midrule
表达形式 & 闭合数学公式 & 近似数值 \\
精度 & 精确无误差 & 有截断与舍入误差 \\
适用范围 & 受限于可求解的级数类型 & 几乎适用于所有收敛级数 \\
参数依赖性（函数项级数） & 可显式表达参数关系 & 需逐点计算 \\
计算效率 & 一次求解永久有效 & 精度越高计算量越大 \\
理论深度 & 揭示内在数学结构 & 偏重数值近似 \\
\bottomrule
\end{tabularx}
\end{table}

值得强调的是，解析解与数值解并非截然对立的关系，而是存在深刻的互补与关联。解析方法可为数值方法提供误差估计和收敛加速的理论依据，数值解可以解析解为基础构造更高效的算法；反之，数值实验的结果有时能为猜想解析形式提供线索。这种互补关系正是本文谱系化框架的理论出发点。
